\chapter{Tuples and Lists}

\textbf{What this chapter gives you:} the two ordered containers you will use all day. A \code{tuple} is a
fixed row of items. A \code{list} is a row you can grow, shrink, and reorder. You will learn to slice them,
sort them, copy them without surprises, and build new ones with comprehensions. A prompt history, a batch of
documents, and a list of embeddings are all lists.

%RELATED: 4, 5, 7, 9, 10, 11

\begin{keyidea}[: sequences, mutable and immutable]
A \emph{sequence} holds items in order and is read by position (index). Strings and bytes are
sequences of one kind of item. Tuples and lists can hold \emph{any} mix of types, even other lists.
\begin{center}
\begin{tabularx}{0.95\linewidth}{@{}l c c l@{}}
\toprule
\textbf{Type} & \textbf{mutable?} & \textbf{hashable?} & \textbf{Literal} \\
\midrule
\code{str}, \code{bytes} & \xmark & \cmark & \code{"ab"}, \code{b"ab"} \\
\code{tuple}             & \xmark & \cmark (if every item is) & \code{(1, 2)} \\
\code{list}              & \cmark & \xmark & \code{[1, 2]} \\
\code{bytearray}         & \cmark & \xmark & \code{bytearray(b"ab")} \\
\bottomrule
\end{tabularx}
\end{center}
\end{keyidea}

% ------------------------------------------------------------
\section{Tuples}

\subsection{Create with commas and ()}

The comma makes the tuple, not the parentheses. Parentheses only make it easier to see. The one place where
the comma really matters is a tuple with one item.

\begin{pycode}
>>> empty = ()
>>> pair = 1024, "tokens"
>>> pair
(1024, 'tokens')
>>> one = ("haiku",)
>>> one, type(one)
(('haiku',), <class 'tuple'>)
>>> not_a_tuple = ("haiku")
>>> type(not_a_tuple)
<class 'str'>
\end{pycode}

\begin{trap}[: the one-item tuple]
\code{("haiku")} is just a string in brackets. You need the trailing comma: \code{("haiku",)}. This also
bites when you pass a one-item tuple to a function: \code{f(("a",))}, because \code{f("a",)} is just the
string.
\end{trap}

\subsection{Unpacking and swapping}

Assigning a tuple to several names takes it apart. The count must match, unless you use \code{*name} to
soak up the extras (it always becomes a list).

\begin{pycode}
>>> usage = (812, 143)
>>> tokens_in, tokens_out = usage
>>> tokens_in, tokens_out
(812, 143)
>>> a, b, c = usage
Traceback (most recent call last):
  ...
ValueError: not enough values to unpack (expected 3, got 2)
>>> first, *rest = [10, 20, 30, 40]
>>> first, rest
(10, [20, 30, 40])
>>> x, y = 1, 2
>>> x, y = y, x
>>> x, y
(2, 1)
\end{pycode}

The swap works because the right side is built into a tuple first, then unpacked. No temporary variable is
needed.

\subsection{Convert with tuple()}

\code{tuple(iterable)} builds a tuple from anything you can loop over.

\begin{pycode}
>>> tuple(["a", "b"]), tuple("ab"), tuple(b"\x01\x02")
(('a', 'b'), ('a', 'b'), (1, 2))
\end{pycode}

\subsection{What you can do with a tuple}

Read it, slice it, combine it, repeat it, compare it, loop over it. You cannot change it.

\begin{pycode}
>>> t = (1, 2, 3)
>>> t[0], t[-1], t[1:]
(1, 3, (2, 3))
>>> t[0] = 9
Traceback (most recent call last):
  ...
TypeError: 'tuple' object does not support item assignment
>>> t + (4,), t * 2
((1, 2, 3, 4), (1, 2, 3, 1, 2, 3))
>>> (7, 2) < (7, 2, 9), (7, 2) == (7, 2, 9)
(True, False)
>>> for word in ("fresh", "out"):
...     print(word)
fresh
out
\end{pycode}

``Modify'' with \code{+=} only \emph{looks} like a change. Python builds a new tuple and moves the name onto
it, exactly as for \code{bytes}.

\begin{pycode}
>>> t1 = ("Fee", "Fie")
>>> before = t1
>>> t1 += ("Foe",)
>>> t1, before, t1 is before
(('Fee', 'Fie', 'Foe'), ('Fee', 'Fie'), False)
\end{pycode}

\begin{trap}[: immutable does not mean ``deeply frozen'']
A tuple cannot re-point its slots, but if a slot holds a \emph{list}, that list can still change. The
tuple is then not hashable.
\end{trap}

\begin{pycode}
>>> record = ("log", [1, 2])
>>> record[1].append(3)
>>> record
('log', [1, 2, 3])
>>> hash(record)
Traceback (most recent call last):
  ...
TypeError: unhashable type: 'list'
\end{pycode}

A \emph{named tuple} lets you read items by name as well as by position. It is covered in Chapter 11.

% ------------------------------------------------------------
\section{Lists}

A list keeps items in order, allows duplicates, and can be changed in place. Use one whenever the contents
or the length can change.

\subsection{Create with [], list(), and split()}

\begin{pycode}
>>> empty = []
>>> mixed = [1, "two", [3, 4], None]
>>> list("cat"), list((1, 2)), list(range(3))
(['c', 'a', 't'], [1, 2], [0, 1, 2])
>>> "9/19/2024".split("/")
['9', '19', '2024']
>>> "a/b//c".split("/")
['a', 'b', '', 'c']
>>> "a/b//c".split("//")
['a/b', 'c']
>>> "one  two   three".split()
['one', 'two', 'three']
\end{pycode}

Two repeated separators give an empty string between them. \code{split()} with no argument splits on any run
of whitespace and drops the empties.

\subsection{Index and slice}

An index counts from 0 at the left, or from $-1$ at the right. A slice \code{seq[start:stop:step]} returns a
\emph{new list} of the same type. The easiest way to read a slice is to imagine the numbers sitting on the
\emph{lines between} items: the slice runs from the start line to the stop line, so the stop item is not
included.

%MD<
% ```text
%  item index:    0     1     2     3     4
%              +-----+-----+-----+-----+-----+
%   row =      |  p  |  q  |  r  |  s  |  t  |
%              +-----+-----+-----+-----+-----+
%  negative:    -5    -4    -3    -2    -1
%  slice lines: 0     1     2     3     4     5
%                     [-------------)           row[1:4] -> [q, r, s]
% ```
%MD>
\begin{center}
\begin{tikzpicture}[cell/.style={draw=cblue,fill=cblue!8,minimum width=12mm,minimum height=8mm,font=\ttfamily\small},
  ix/.style={font=\ttfamily\scriptsize,text=cgray},ln/.style={font=\ttfamily\scriptsize,text=cgold!80!black},
  lab/.style={font=\sffamily\scriptsize,text=cgray},>=Stealth]
  \foreach \i/\v in {0/p,1/q,2/r,3/s,4/t} {
    \node[cell] at (1.2*\i,0) {\v};
    \node[ix] at (1.2*\i,0.7) {\i};
  }
  \foreach \i/\n in {0/-5,1/-4,2/-3,3/-2,4/-1} \node[ix] at (1.2*\i,-0.7) {\n};
  \foreach \k in {0,...,5} {
    \draw[cgold,thick] (1.2*\k-0.6,-0.4) -- (1.2*\k-0.6,-1.1);
    \node[ln] at (1.2*\k-0.6,-1.35) {\k};
  }
  \fill[cgreen,opacity=0.18] (0.6,-0.4) rectangle (4.2,0.4);
  \draw[cgreen!70!black,very thick,|-|] (0.6,-1.8) -- node[below,font=\ttfamily\scriptsize]{row[1:4] gives [q, r, s]} (4.2,-1.8);
  \node[lab,anchor=east] at (-0.9,0.7) {item index};
  \node[lab,anchor=east] at (-0.9,-0.7) {negative index};
  \node[lab,anchor=east] at (-0.9,-1.35) {slice lines};
\end{tikzpicture}
\end{center}

\begin{pycode}
>>> row = ["p", "q", "r", "s", "t"]
>>> row[0], row[-1]
('p', 't')
>>> row[1:4], row[:2], row[-2:]
(['q', 'r', 's'], ['p', 'q'], ['s', 't'])
>>> row[::2], row[::-1]
(['p', 'r', 't'], ['t', 's', 'r', 'q', 'p'])
>>> row[10:], row[-9:2]
([], ['p', 'q'])
>>> row[5]
Traceback (most recent call last):
  ...
IndexError: list index out of range
\end{pycode}

\begin{typecheck}[: index vs slice]
\code{row[2]} returns \emph{one item} (any type; here a \code{str}) and raises \code{IndexError} when out of
range. \code{row[2:3]} returns a \emph{list} (here with one item) and never raises: out-of-range ends
quietly snap to the edge. A slice is a copy, so assigning to it does not change \code{row}.
\end{typecheck}

\subsection{Add items}

\begin{pycode}
>>> q = ["a", "b", "c"]
>>> q.append("d")
>>> q.insert(0, "z")
>>> q.extend(["x", "y"])
>>> q
['z', 'a', 'b', 'c', 'd', 'x', 'y']
>>> q.append(["m", "n"])
>>> q[-1]
['m', 'n']
>>> q += ["end"]
>>> print(q.append("e"))
None
\end{pycode}

\begin{trap}[: append vs extend]
\code{append(x)} adds \emph{one item}, even if \code{x} is a list. \code{extend(it)} (and \code{+=})
adds each item of \code{it}. Using \code{append} with a list gives a list inside the list.
\end{trap}

\begin{trap}[: methods that change a list return None]
\code{append}, \code{extend}, \code{insert}, \code{sort}, \code{reverse}, \code{remove}, and \code{clear}
change the list in place and return \code{None}. Writing \code{names = names.sort()} replaces your list with
\code{None}.
\end{trap}

\subsection{Change items}

Assigning to an index replaces one item. Assigning to a slice replaces a whole stretch, and the new stretch
can be any iterable of any length.

\begin{pycode}
>>> n = [1, 2, 3, 4]
>>> n[0] = "one"
>>> n[1:3] = [7, 8, 9]
>>> n
['one', 7, 8, 9, 4]
>>> n[1:3] = []
>>> n
['one', 9, 4]
>>> n[10] = 1
Traceback (most recent call last):
  ...
IndexError: list assignment index out of range
\end{pycode}

\subsection{Remove items}

\begin{center}
\begin{tabularx}{\linewidth}{@{}l l l X@{}}
\toprule
\textbf{Want to} & \textbf{Use} & \textbf{Returns} & \textbf{If missing} \\
\midrule
delete by position   & \code{del m[i]}   & (statement)  & \code{IndexError} \\
delete by value      & \code{m.remove(v)} & \code{None} & \code{ValueError}; removes the first match only \\
take out by position & \code{m.pop(i)}   & the item     & \code{IndexError}; default \code{i} is \code{-1} \\
empty the list       & \code{m.clear()}  & \code{None}  & -- \\
\bottomrule
\end{tabularx}
\end{center}

\begin{pycode}
>>> m = ["a", "b", "c", "b"]
>>> del m[0]
>>> m.remove("b")
>>> m
['c', 'b']
>>> m.pop(), m.pop(0)
('b', 'c')
>>> m.pop()
Traceback (most recent call last):
  ...
IndexError: pop from empty list
>>> ["a"].remove("z")
Traceback (most recent call last):
  ...
ValueError: list.remove(x): x not in list
\end{pycode}

\code{append()} with \code{pop()} gives a stack (last in, first out). \code{append()} with \code{pop(0)} gives
a queue (first in, first out), but \code{pop(0)} is slow on big lists. For queues, use
\code{collections.deque} (Chapter 12).

\subsection{Find and count}

\begin{pycode}
>>> w = ["a", "b", "a"]
>>> w.index("a"), w.count("a"), "b" in w, len(w)
(0, 2, True, 3)
>>> w.index("z")
Traceback (most recent call last):
  ...
ValueError: 'z' is not in list
>>> nums = [3, 9, 4]
>>> min(nums), max(nums), sum(nums)
(3, 9, 16)
\end{pycode}

\code{index()} returns the \emph{first} position and raises \code{ValueError} if the value is absent, so check
with \code{in} first when it may be missing. If you test membership a lot, a set is faster (Chapter 9).

\subsection{join() is the opposite of split()}

\code{join()} is a \emph{string} method. The separator is the string you call it on, and the list is the
argument. It returns a \code{str} and needs every item to already be a string.

\begin{pycode}
>>> parts = ["Harry", "Hermione", "Ron"]
>>> ", ".join(parts)
'Harry, Hermione, Ron'
>>> ", ".join(parts).split(", ") == parts
True
>>> ", ".join([1, 2])
Traceback (most recent call last):
  ...
TypeError: sequence item 0: expected str instance, int found
>>> ", ".join([str(n) for n in [1, 2]])
'1, 2'
\end{pycode}

\subsection{Sort with sort() or sorted()}

\begin{center}
\begin{tabularx}{\linewidth}{@{}l l l X@{}}
\toprule
 & \textbf{Changes the list?} & \textbf{Returns} & \textbf{Works on} \\
\midrule
\code{m.sort()}   & yes, in place & \code{None} & lists only \\
\code{sorted(x)}  & no            & a new \code{list} & any iterable \\
\code{m.reverse()} & yes, in place & \code{None} & lists only \\
\code{reversed(x)} & no            & an iterator & any sequence \\
\bottomrule
\end{tabularx}
\end{center}

\begin{pycode}
>>> scores = [0.9, 0.2, 0.5]
>>> sorted(scores)
[0.2, 0.5, 0.9]
>>> scores
[0.9, 0.2, 0.5]
>>> result = scores.sort(reverse=True)
>>> result is None, scores
(True, [0.9, 0.5, 0.2])
>>> sorted([1, "a"])
Traceback (most recent call last):
  ...
TypeError: '<' not supported between instances of 'str' and 'int'
\end{pycode}

The \code{key=} argument names a function that turns each item into the value to compare. The items
themselves are not changed.

\begin{pycode}
>>> words = ["Banana", "apple", "cherry"]
>>> sorted(words)
['Banana', 'apple', 'cherry']
>>> sorted(words, key=str.lower)
['apple', 'Banana', 'cherry']
>>> runs = [("haiku", 0.4), ("opus", 3.1), ("sonnet", 1.2)]
>>> sorted(runs, key=lambda r: r[1], reverse=True)
[('opus', 3.1), ('sonnet', 1.2), ('haiku', 0.4)]
\end{pycode}

(\code{lambda r: r[1]} is a tiny unnamed function, which Chapter 10 explains. It means ``given \code{r},
give back \code{r[1]}''.) Plain uppercase letters sort before lowercase, and a sort that is
\emph{stable} keeps the original order of ties.

\subsection{Assignment makes a second name, not a second list}

A variable is a label. \code{b = a} sticks a second label on the \emph{same} list, so a change through one
name shows through the other. To get an independent list, copy it with \code{a.copy()}, \code{list(a)}, or
\code{a[:]}.

%MD<
% ```text
%   b = a                              b = a.copy()
%
%   a ----\                            a ---> [1, 2, 3]
%          +--> [1, 2, 3]
%   b ----/                            b ---> [1, 2, 3]     (a second list)
%
%   b is a  -> True                    b is a  -> False
%   a[0] = 9  also changes b           a[0] = 9  leaves b alone
% ```
%MD>
\begin{center}
\begin{tikzpicture}[box/.style={draw=cblue,fill=cblue!8,minimum height=8mm,minimum width=18mm,font=\ttfamily\small},
  nm/.style={draw=cgold,fill=cgold!15,rounded corners=2pt,font=\ttfamily\small,inner sep=3pt},
  lab/.style={font=\sffamily\scriptsize,text=cgray,align=center},>=Stealth]
  \node[lab] at (1.8,2.1) {\textbf{\code{b = a}}};
  \node[nm] (a1) at (0,1.3) {a};
  \node[nm] (b1) at (0,0.2) {b};
  \node[box] (l1) at (3.0,0.75) {[1, 2, 3]};
  \draw[->,thick] (a1) -- (l1.west);
  \draw[->,thick] (b1) -- (l1.west);
  \node[lab] at (1.8,-0.7) {one list, two names\\\code{b is a} is \code{True}};
  \node[lab] at (8.4,2.1) {\textbf{\code{b = a.copy()}}};
  \node[nm] (a2) at (6.4,1.3) {a};
  \node[nm] (b2) at (6.4,0.2) {b};
  \node[box] (l2) at (9.2,1.3) {[1, 2, 3]};
  \node[box] (l3) at (9.2,0.2) {[1, 2, 3]};
  \draw[->,thick] (a2) -- (l2);
  \draw[->,thick] (b2) -- (l3);
  \node[lab] at (8.4,-0.7) {two lists, equal contents\\\code{b is a} is \code{False}};
\end{tikzpicture}
\end{center}

\begin{pycode}
>>> a = [1, 2, 3]
>>> b = a
>>> a[0] = "changed"
>>> b
['changed', 2, 3]
>>> b is a
True
>>> c, d, e = a.copy(), list(a), a[:]
>>> a[1] = "x"
>>> c, d == c, e is a
(['changed', 2, 3], True, False)
\end{pycode}

\subsection{Copy everything with deepcopy()}

\code{copy()}, \code{list()}, and \code{[:]} are \emph{shallow}: they copy the outer list but reuse the inner
objects. If an item is itself a list, both outer lists point at the \emph{same} inner list. For a fully
independent copy, use \code{copy.deepcopy()} from the standard library.

%MD<
% ```text
%   a = [1, 2, X] ----.
%   b = a.copy()       +--> X = [8, 10]    (shared inner list, shallow)
%         [1, 2, X] --'
%
%   c = deepcopy(a)  [1, 2, Y] ---> Y = [8, 9]   (its own inner list)
% ```
%MD>
\begin{center}
\begin{tikzpicture}[cell/.style={draw=cblue,fill=cblue!8,minimum width=8mm,minimum height=7mm,font=\ttfamily\small},
  nm/.style={draw=cgold,fill=cgold!15,rounded corners=2pt,font=\ttfamily\small,inner sep=3pt},
  inner/.style={draw=cred,fill=cred!8,minimum height=7mm,font=\ttfamily\small,inner sep=4pt},
  lab/.style={font=\sffamily\scriptsize,text=cgray,align=center},>=Stealth]
  % a
  \node[nm] at (-0.9,2.2) {a};
  \node[cell] at (0,2.2) {1};  \node[cell] at (0.8,2.2) {2};  \node[cell] (oa) at (1.6,2.2) {$\bullet$};
  % b
  \node[nm] at (3.3,2.2) {b};
  \node[cell] at (4.2,2.2) {1};  \node[cell] at (5.0,2.2) {2};  \node[cell] (ob) at (5.8,2.2) {$\bullet$};
  % c
  \node[nm] at (8.5,2.2) {c};
  \node[cell] at (9.4,2.2) {1};  \node[cell] at (10.2,2.2) {2};  \node[cell] (oc) at (11.0,2.2) {$\bullet$};
  \node[inner] (shared) at (3.7,0.3) {[8, 10]};
  \node[inner] (own) at (11.0,0.3) {[8, 9]};
  \draw[->,thick,cred] (oa.south) -- (shared.north west);
  \draw[->,thick,cred] (ob.south) -- (shared.north east);
  \draw[->,thick,cred] (oc.south) -- (own.north);
  \node[lab] at (0.8,3.0) {original};
  \node[lab] at (5.0,3.0) {\code{b = a.copy()}};
  \node[lab] at (10.2,3.0) {\code{c = deepcopy(a)}};
  \node[lab,anchor=west] at (0,-0.6) {after \code{a[2][1] = 10}: \code{a} and \code{b} see \code{[8, 10]}; \code{c} still has \code{[8, 9]}};
\end{tikzpicture}
\end{center}

\begin{pycode}
>>> import copy
>>> a = [1, 2, [8, 9]]
>>> b = a.copy()
>>> c = copy.deepcopy(a)
>>> a[2][1] = 10
>>> b, c
([1, 2, [8, 10]], [1, 2, [8, 9]])
>>> b[2] is a[2], c[2] is a[2]
(True, False)
\end{pycode}

\begin{boss}[: two names, one list (and shared insides)]
Almost every strange ``my data changed by itself'' bug is aliasing. \code{b = a} shares the list. Even
\code{a.copy()} shares any \emph{inner} lists, dicts, or sets. This applies to \emph{every} mutable type.\\[2pt]
\textbf{Strategy:} ask ``is this one object or two?'' with \code{is}. Copy when you need independence. Use
\code{deepcopy} when the data is nested. Prefer building new lists over editing a shared one.
\end{boss}

\begin{trap}[: multiplying a list of lists repeats one inner list]
Multiplying a list repeats \emph{references}, so \code{[[0] * 3] * 2} holds the same inner list twice. The rows of the grid below are one list shown twice. Build
each row separately with a comprehension.
\end{trap}

\begin{pycode}
>>> grid = [[0] * 3] * 2
>>> grid[0][0] = 1
>>> grid
[[1, 0, 0], [1, 0, 0]]
>>> grid = [[0] * 3 for _ in range(2)]
>>> grid[0][0] = 1
>>> grid
[[1, 0, 0], [0, 0, 0]]
\end{pycode}

\subsection{Compare and iterate}

Lists are compared item by item from the left, like words in a dictionary. A shorter list that matches
the start is smaller. Iteration is \code{for item in list}; \code{break}, \code{continue}, and \code{else}
work as in Chapter 7.

\begin{pycode}
>>> [7, 2] == [7, 2, 9], [7, 2] < [7, 2, 9], [3] < [2, 9]
(False, True, False)
>>> for cheese in ["brie", "gjetost", "havarti"]:
...     if cheese.startswith("g"):
...         print("skip", cheese)
...         break
...     print(cheese)
... else:
...     print("no g found")
brie
skip gjetost
\end{pycode}

\subsection{Iterate several sequences with zip()}

\code{zip(a, b, ...)} pairs up items at the same position and stops at the shortest input. It returns a
lazy iterator of tuples, so wrap it in \code{list()} to see them, and you can unpack each tuple in the
\code{for} line.

%MD<
% ```text
%   models  :  haiku     sonnet    opus
%   prices  :  0.1       0.3       1.0
%              |         |         |
%   zip ->  (haiku,0.1) (sonnet,0.3) (opus,1.0)     one tuple per position
% ```
%MD>
\begin{center}
\begin{tikzpicture}[c/.style={draw=cblue,fill=cblue!8,minimum width=16mm,minimum height=7mm,font=\ttfamily\small},
  p/.style={draw=cgreen,fill=cgreen!10,minimum width=16mm,minimum height=7mm,font=\ttfamily\small},
  lab/.style={font=\sffamily\scriptsize,text=cgray},>=Stealth]
  \foreach \i/\m/\pr in {0/haiku/0.1,1/sonnet/0.3,2/opus/1.0} {
    \node[c] (m\i) at (2.4*\i,1.0) {\m};
    \node[p] (p\i) at (2.4*\i,0) {\pr};
    \node[draw=cpurple,rounded corners=3pt,thick,fit=(m\i)(p\i),inner sep=2.5pt] {};
  }
  \node[lab,anchor=east] at (-1.2,1.0) {\code{models}};
  \node[lab,anchor=east] at (-1.2,0) {\code{prices}};
  \node[lab,anchor=west] at (6.2,0.5) {each purple box is one tuple};
\end{tikzpicture}
\end{center}

\begin{pycode}
>>> models = ["haiku", "sonnet", "opus"]
>>> prices = [0.1, 0.3, 1.0]
>>> for name, price in zip(models, prices):
...     print(name, price)
haiku 0.1
sonnet 0.3
opus 1.0
>>> dict(zip(models, prices))
{'haiku': 0.1, 'sonnet': 0.3, 'opus': 1.0}
>>> list(zip("ab", [1, 2, 3]))
[('a', 1), ('b', 2)]
\end{pycode}

If the lengths differ, \code{zip()} silently drops the extras. Two ways to stay in control:

\begin{pycode}
>>> from itertools import zip_longest
>>> list(zip_longest([1, 2], ["a", "b", "c"], fillvalue=0))
[(1, 'a'), (2, 'b'), (0, 'c')]
>>> list(zip([1], [1, 2], strict=True))
Traceback (most recent call last):
  ...
ValueError: zip() argument 2 is longer than argument 1
\end{pycode}

\subsection{List comprehensions}

A comprehension builds a list in one expression: \code{[expression for item in iterable if condition]}.
The \code{if} part is optional. It reads as: for each item, if the condition is true, put the expression's
value in the new list. The loop moves inside the brackets.

%MD<
% ```text
%   [  n * 2        for n in range(10)       if n % 2 == 0   ]
%      ^ (3)           ^ (1)                    ^ (2)
%      value to add    take each item           keep only if True
% ```
%MD>
\begin{center}
\begin{tikzpicture}[chip/.style={draw=#1,fill=#1!10,rounded corners=2pt,font=\ttfamily\small,inner sep=4pt,minimum height=7mm},
  num/.style={circle,fill=cgold,text=white,font=\sffamily\scriptsize\bfseries,inner sep=1.5pt},
  lab/.style={font=\sffamily\scriptsize,text=cgray,align=center},>=Stealth]
  \node[font=\ttfamily\large] (lb) at (-0.4,0) {[};
  \node[chip=cgreen,right=1mm of lb] (e) {n * 2};
  \node[chip=cblue,right=3mm of e] (f) {for n in range(10)};
  \node[chip=cgold,right=3mm of f] (g) {if n \% 2 == 0};
  \node[font=\ttfamily\large,right=1mm of g] (rb) {]};
  \node[num,above=3mm of f] {1};
  \node[num,above=3mm of g] {2};
  \node[num,above=3mm of e] {3};
  \node[lab,below=3mm of e] {value that\\goes into the list};
  \node[lab,below=3mm of f] {take each item\\from the iterable};
  \node[lab,below=3mm of g] {keep it only\\if this is True};
\end{tikzpicture}
\end{center}

\begin{pycode}
>>> [n * 2 for n in range(4)]
[0, 2, 4, 6]
>>> [n for n in range(10) if n % 3 == 0]
[0, 3, 6, 9]
>>> words = ["a", "bb", "ccc"]
>>> [len(w) for w in words]
[1, 2, 3]
>>> [(r, c) for r in range(2) for c in range(2)]
[(0, 0), (0, 1), (1, 0), (1, 1)]
\end{pycode}

With two \code{for} clauses, the first is the outer loop. The comprehension above equals a nested loop that
appends \code{(r, c)} at the bottom.

\begin{trap}[: an \texttt{if} at the end filters; an \texttt{if/else} at the front transforms]
These look alike but do different jobs. A trailing \code{if} (no \code{else}) \emph{drops} items. A
conditional expression in the value slot keeps every item and picks a value for each. Remember the order
from Chapter 6: value, \code{if} test, \code{else} value.
\end{trap}

\begin{pycode}
>>> words = ["a", "bb", "ccc"]
>>> [w for w in words if len(w) > 1]
['bb', 'ccc']
>>> ["long" if len(w) > 1 else "short" for w in words]
['short', 'long', 'long']
\end{pycode}

\subsection{Lists of lists}

A list can hold lists. Use one index per level, left to right: the first picks the inner list, the next
picks inside it.

\begin{pycode}
>>> batches = [["a", "b"], ["c"], []]
>>> batches[0], batches[0][1]
(['a', 'b'], 'b')
>>> len(batches), [len(b) for b in batches]
(3, [2, 1, 0])
>>> [item for batch in batches for item in batch]
['a', 'b', 'c']
\end{pycode}

The last line \emph{flattens} one level. The outer \code{for} comes first, just as in nested loops.

% ------------------------------------------------------------
\section{Tuples versus lists}

You can often use either. The choice is about what you mean.

\begin{center}
\begin{tabularx}{\linewidth}{@{}l X X@{}}
\toprule
 & \textbf{tuple} & \textbf{list} \\
\midrule
Mutable?              & no                                       & yes \\
Methods               & only \code{count} and \code{index}        & many (\code{append}, \code{sort}, ...) \\
Can be a dict key?    & yes, if all items are hashable           & no \\
Typically holds       & a fixed record of \emph{different} things: \code{(model, price)} & many \emph{similar} items: a list of prices \\
Size                  & a little smaller                         & over-allocates room to grow \\
Safe from accidents?  & yes, nobody can edit it by mistake       & no \\
\bottomrule
\end{tabularx}
\end{center}

\begin{gamerule}
Use a tuple for a small fixed record or a function's multiple results. Use a list for a collection that
grows, shrinks, or gets sorted. In everyday code you will use lists (and dicts, Chapter 9) far more.
\end{gamerule}

% ------------------------------------------------------------
\section{There are no tuple comprehensions}

Lists, dicts, and sets have comprehensions. Tuples, strings, and bytes do not. Swapping the square brackets
for parentheses does \emph{not} give a tuple. It gives a \emph{generator expression}: a lazy iterator that
produces items one at a time and can be read only once. Wrap it in \code{tuple()} to get a tuple.

\begin{pycode}
>>> g = (n * n for n in range(3))
>>> type(g)
<class 'generator'>
>>> tuple(g)
(0, 1, 4)
>>> tuple(g)
()
>>> sum(n * n for n in range(4))
14
\end{pycode}

Because a generator makes values only when asked, it is the right tool for big or endless data, for example
processing a million documents without holding them all in memory. Chapter 10 covers generators fully.

\begin{practice}
\begin{enumerate}
\item Make a list \code{prices = [0.1, 0.3, 1.0, 0.6]}. Write the slices that give the first two, the last
two, and the list reversed. What is each result's type?
\item What do \code{x = [3, 1, 2].sort()} and \code{y = sorted([3, 1, 2])} leave in \code{x} and \code{y}?
\item After \code{a = [1, [2, 3]]; b = a.copy(); b[1].append(4)}, what is \code{a}? Why? What one change makes
\code{a} stay untouched?
\item Write a list comprehension that gives the lengths of only the words longer than 3 letters in a list of
words.
\item What is the difference between \code{append} and \code{extend}? Show each with \code{[1, 2]} and
\code{[3, 4]}.
\item Why does \code{("hello")} not make a tuple, and what makes the one-item tuple?
\end{enumerate}
\end{practice}

\begin{achievement}{List Keeper}
You can now:
\begin{itemize}
\item build tuples and lists, unpack them, and slice with start, stop, and step;
\item add, change, find, and remove items, knowing which methods return \code{None};
\item sort with \code{sort}, \code{sorted}, and \code{key=}, and pair sequences with \code{zip};
\item tell aliasing, shallow copy, and deep copy apart and pick the right one;
\item write list comprehensions with filters and conditional values.
\end{itemize}
\end{achievement}
